% ============================================================
% PHẦN A — KIẾN THỨC NỀN GRAPHVIS — SƯỜN
% Chỉ trình bày những gì GraphVis thực sự dùng.
% ============================================================
\section{GraphVis: kiến thức nền}\label{sec:gv-background}

\subsection{Đồ thị tri thức và ConceptNet}
\todo{Định nghĩa KG $\mathcal{G}=(\mathcal{V},\mathcal{E})$ với bộ ba (đầu, quan hệ, đuôi); KG thường thức vs KG thực thể; ConceptNet 5.5: nút là khái niệm đời thường, cạnh có hướng và nhãn; bản tiếng Anh dùng trong GraphVis có 799.273 nút, 2.487.810 cạnh, 42 quan hệ gộp còn 17.}

\subsection{Mô hình ngôn ngữ sinh và tinh chỉnh có giám sát}
\todo{Phân rã tự hồi quy $p_\theta(\mathbf{y}\mid\mathbf{x})=\prod_t p_\theta(y_t\mid\mathbf{x},\mathbf{y}_{<t})$; hàm mất mát NLL của SFT (sẽ dùng lại ở \eqref{eq:gv-loss}); giải mã tham lam khi đánh giá.}

\subsection{Mô hình ngôn ngữ--thị giác lớn (LVLM)}
\todo{Cấu trúc LLaVA: bộ mã hóa ảnh CLIP ViT-L/14 $f_v$, bộ chiếu $f_p$, LLM; ảnh được biến thành chuỗi token thị giác ghép trước token văn bản; LLaVA-1.6 (LLaVA-NeXT) dùng AnyRes khi suy luận. Hình nhỏ vẽ luồng ảnh $\rightarrow$ token.}

\subsection{LoRA}
\todo{$W=W_0+\frac{\alpha}{r}BA$, $B\in\mathbb{R}^{d\times r}$, $A\in\mathbb{R}^{r\times k}$, $r\ll\min(d,k)$; chỉ cập nhật $A,B$; QLoRA (lượng tử hóa 4-bit $W_0$) — phương án nhóm dùng khi tái lập trên Colab nếu bộ nhớ không đủ.}

\subsection{Học theo giáo trình (curriculum learning)}
\todo{Ý tưởng của Bengio và cộng sự: sắp xếp dữ liệu/tác vụ từ dễ đến khó; trong GraphVis ``giáo trình'' ở mức tác vụ (đọc hình $\rightarrow$ trả lời câu hỏi), không phải mức mẫu.}

